% BEGIN VOL07-EXPANSION VII27
\section{Supplementary worked examples}
\label{sec:vii27-pedagogy}

\begin{example}[Round sphere contractions]\label{ex:vii27-ped-01}
On \(S^n(a)\), constant sectional curvature \(1/a^2\) gives \(\operatorname{Ric}=\frac{n-1}{a^2}g\) and \(\operatorname{Scal}=\frac{n(n-1)}{a^2}\).  The contractions retain the radius scale while discarding directional detail.
\end{example}

\begin{example}[Hyperbolic space]\label{ex:vii27-ped-02}
On hyperbolic \(n\)-space of curvature \(-1/a^2\), \(\operatorname{Ric}=-(n-1)a^{-2}g\) and \(\operatorname{Scal}=-n(n-1)a^{-2}\).  Negative sectional curvature therefore gives negative Ricci in every direction.
\end{example}

\begin{example}[Product \(S^2(a)\times\mathbb R\)]\label{ex:vii27-ped-03}
In the product metric on \(S^2(a)\times\mathbb R\), Ricci eigenvalues are \(1/a^2,1/a^2,0\).  The scalar curvature is \(2/a^2\): the flat factor contributes nothing, while the spherical factor contributes both curved directions.
\end{example}

\section{Graded supplementary exercises}
The following exercises add a deliberate progression from concrete calculation to structural proof, hypothesis testing, application, and synthesis.

\subsection*{Standard computations and constructions}

\begin{exercise}[Ricci sum]\label{exr:vii27-ped-01}
For a unit vector \(v\), how is \(\operatorname{Ric}(v,v)\) expressed using sectional curvatures?
\end{exercise}

\begin{hint}
Complete \(v\) to an orthonormal basis.
\end{hint}

\begin{solution}
It is the sum of sectional curvatures of the \(n-1\) planes containing \(v\).
\end{solution}

\begin{exercise}[Scalar trace]\label{exr:vii27-ped-02}
How is scalar curvature obtained from Ricci?
\end{exercise}

\begin{hint}
Take the metric trace.
\end{hint}

\begin{solution}
\(\operatorname{Scal}=g^{ij}\operatorname{Ric}_{ij}\).
\end{solution}

\begin{exercise}[Space form]\label{exr:vii27-ped-03}
In dimension \(4\) with constant sectional curvature \(c\), compute scalar curvature.
\end{exercise}

\begin{hint}
Use \(n(n-1)c\).
\end{hint}

\begin{solution}
\(\operatorname{Scal}=12c\).
\end{solution}

\begin{exercise}[Sphere two]\label{exr:vii27-ped-04}
Compute scalar curvature of \(S^2(a)\).
\end{exercise}

\begin{hint}
Use \(n=2\).
\end{hint}

\begin{solution}
\(\operatorname{Scal}=2/a^2\).
\end{solution}

\begin{exercise}[Flat space]\label{exr:vii27-ped-05}
What are Ricci and scalar curvature in Euclidean space?
\end{exercise}

\begin{hint}
All sectional curvatures vanish.
\end{hint}

\begin{solution}
Both are zero.
\end{solution}

\subsection*{Proofs}

\begin{exercise}[Ricci symmetry]\label{exr:vii27-ped-06}
Why is the Levi--Civita Ricci tensor symmetric?
\end{exercise}

\begin{hint}
Use Riemann symmetries and Bianchi.
\end{hint}

\begin{solution}
The contraction of the Riemann tensor combined with its pair symmetries and first Bianchi identity yields \(\operatorname{Ric}(X,Y)=\operatorname{Ric}(Y,X)\).
\end{solution}

\begin{exercise}[Scalar sectional sum]\label{exr:vii27-ped-07}
Prove scalar curvature is twice the sum of sectional curvatures over orthonormal basis planes.
\end{exercise}

\begin{hint}
Expand the trace of Ricci.
\end{hint}

\begin{solution}
Each unordered pair \(i<j\) appears twice in \(\sum_i\operatorname{Ric}(e_i,e_i)\), giving the factor two.
\end{solution}

\begin{exercise}[Einstein trace]\label{exr:vii27-ped-08}
If \(\operatorname{Ric}=\lambda g\) in dimension \(n\), compute scalar curvature.
\end{exercise}

\begin{hint}
Take the metric trace.
\end{hint}

\begin{solution}
\(\operatorname{Scal}=n\lambda\).
\end{solution}

\begin{exercise}[Product scalar]\label{exr:vii27-ped-09}
Explain why scalar curvature of a Riemannian product is the sum of scalar curvatures of the factors.
\end{exercise}

\begin{hint}
Use an orthonormal product basis and mixed sectional curvature zero.
\end{hint}

\begin{solution}
The trace splits into factor directions; mixed planes contribute zero, so the scalar curvatures add.
\end{solution}

\subsection*{Counterexamples and hypothesis tests}

\begin{exercise}[Ricci zero implies flat]\label{exr:vii27-ped-10}
Test the claim: Ricci-flat always means Riemann-flat.
\end{exercise}

\begin{hint}
Think of dimensions at least four.
\end{hint}

\begin{solution}
False.  Weyl curvature can remain nonzero even when Ricci vanishes.
\end{solution}

\begin{exercise}[Scalar zero implies flat]\label{exr:vii27-ped-11}
Test the claim: zero scalar curvature implies zero Ricci curvature.
\end{exercise}

\begin{hint}
A trace can vanish while a tensor does not.
\end{hint}

\begin{solution}
False.  Positive and negative Ricci eigenvalues can cancel in the trace.
\end{solution}

\begin{exercise}[Positive scalar all planes]\label{exr:vii27-ped-12}
Test the claim: positive scalar curvature means every sectional curvature is positive.
\end{exercise}

\begin{hint}
Scalar curvature is an average sum.
\end{hint}

\begin{solution}
False.  Some planes may have negative curvature while the total sum remains positive.
\end{solution}

\subsection*{Applications and investigations}

\begin{exercise}[Volume distortion]\label{exr:vii27-ped-13}
Why is Ricci curvature central in comparison geometry?
\end{exercise}

\begin{hint}
Think about averages of sectional curvature around one direction.
\end{hint}

\begin{solution}
It controls averaged transverse focusing and appears in volume-growth and geodesic-congruence comparison estimates.
\end{solution}

\begin{exercise}[Einstein equation bridge]\label{exr:vii27-ped-14}
Why is Ricci, rather than the full Riemann tensor, prominent in gravitational field equations?
\end{exercise}

\begin{hint}
It is the canonical two-tensor contraction of curvature.
\end{hint}

\begin{solution}
It has the same tensor type as the metric and stress-energy tensor after forming the Einstein tensor, making a covariant rank-two field equation possible.
\end{solution}

\subsection*{Challenge problems}

\begin{exercise}[Product eigenvalues]\label{exr:vii27-ped-15}
Verify the Ricci eigenvalues of \(S^2(a)\times\mathbb R\).
\end{exercise}

\begin{hint}
Use product splitting.
\end{hint}

\begin{solution}
The sphere directions each have Ricci \(1/a^2\); the line direction has zero Ricci.
\end{solution}

\begin{exercise}[Recover curvature in dimension two]\label{exr:vii27-ped-16}
Why does scalar curvature determine the full Riemann tensor on a surface?
\end{exercise}

\begin{hint}
There is only one tangent two-plane.
\end{hint}

\begin{solution}
In dimension two \(\operatorname{Scal}=2K\), and the Riemann tensor is determined by the single Gaussian curvature \(K\).
\end{solution}

% END VOL07-EXPANSION VII27
